
A partir do desenvolvimento do projeto LibrasVision, conclui-se que a aplicação de técnicas de Visão Computacional e Aprendizado Profundo apresenta grande potencial para promover a inclusão de pessoas surdas no ambiente acadêmico. A proposta de utilizar modelos de Pose Estimation aliados a redes neurais recorrentes mostrou-se adequada para o reconhecimento de sinais dinâmicos, considerando sua capacidade de interpretar movimentos ao longo do tempo.

Além disso, o projeto evidencia que soluções tecnológicas podem reduzir barreiras comunicacionais, proporcionando maior autonomia aos estudantes surdos e diminuindo a dependência de intérpretes em interações cotidianas. A utilização do MediaPipe para extração de landmarks contribui significativamente para a eficiência do sistema, ao reduzir o custo computacional e preservar a privacidade dos usuários.

Por fim, espera-se que o protótipo alcance níveis satisfatórios de precisão e baixa latência, tornando viável sua aplicação em tempo real. Como continuidade, o projeto pode ser expandido com o aumento do vocabulário, aprimoramento dos modelos e testes em cenários reais, consolidando-se como uma ferramenta relevante no avanço da acessibilidade e da interação humano-computador.

\section{Trabalhos futuros}

